% ============================================================
% showcase-report.tex — harryopo-report.cls 全功能展示模板
% 覆盖：封面/摘要/目录/引言、章-节-小节-细目层级、
%        中英文混排、公式、表格、图片、代码、定理、引用
% 编译：xelatex(3次)
% ============================================================
\documentclass{harryopo-report}

% --- 报告信息 ---
\title{harryopo LaTeX 报告模板\\全功能展示手册}
\subtitle{——从封面到参考文献的完整排版指南}
\author{模板作者\quad 测试用户}
\date{2026年6月21日}

\begin{document}

% ==== 封面 ====
\maketitle

% ==== 摘要 ====
\begin{abstract}
本文档是学术报告模板（\texttt{harryopo-report.cls}）的全功能展示手册，演示了从封面、摘要、目录、引言到正文各章的全流程排版能力。文档基于 \texttt{ctexrep} 文档类构建，提供章、节、小节、小小节、括号细目五级标题层级；集成 \texttt{unicode-math} + XITS Math 数学引擎；支持三线表、自适应列宽表、多行合并表等多种表格样式；提供 TikZ 内联绘图与代码高亮（\texttt{listings}）功能。读者可直接复制本文档中的代码片段进行二次开发，快速构建规范、专业的学术报告。
\end{abstract}

% ==== 目录 ====
\tableofcontents

% ==== 引言 ====
\introduction

学术报告是高校教学和科研活动中常见的文档类型。一份排版规范、结构清晰的报告不仅能提升阅读体验，还能体现作者的专业素养。然而，许多初学者在使用 LaTeX 时面临模板配置复杂、字体设置繁琐、中英文混排困难等问题。

本模板就是为了解决这些痛点而设计的。它封装了全部排版细节——从页边距、字体、颜色到章节格式、图表题注、参考文献样式——用户只需专注于内容写作，无需关心排版实现。以下各章将逐一展示模板的全部功能特性。

% ============================================================
\chapter{文字与段落排版}

\section{中英文混排}

中文段落中可自然嵌入英文词汇、缩写和数字。例如，BERT（Bidirectional Encoder Representations from Transformers）采用自监督预训练（Self-Supervised Pre-training）策略，在 11 项 NLP 基准测试上达到了当时的 state-of-the-art 性能。XITS 字体族确保中英文混排时视觉风格统一，数字如 $3.14159$ 与英文如 \emph{overfitting} 与中文如"过拟合"之间过渡自然。

The development of Large Language Models (LLMs) has been one of the most significant technological breakthroughs in recent years. Models such as GPT-4, LLaMA 3, and DeepSeek-V3 have demonstrated remarkable capabilities in reasoning, coding, and multilingual understanding. 这些模型通常包含数千亿个参数，训练的算力成本高达数百万美元。

\section{多级标题层级展示}

报告模板提供五级标题：章（\verb|\chapter|）、节（\verb|\section|）、小节（\verb|\subsection|）、小小节（\verb|\subsubsection|）、括号细目（\verb|\subhead|）。

\subsection{小节标题（黑体 + 辅色）}
小节的标题字体为黑体加粗，颜色为辅助主题色，左缩进 1em。

\subsubsection{小小节标题（楷体）}
小小节使用方正楷体，左缩进 2em。适用于更细粒度的主题拆分。

\subhead{括号细目（无编号）}
括号式标题不产生编号，以黑体标注于行首，适合逻辑分割但无需结构层级编号的场景。

\subhead{另一个括号细目}通过 \verb|\subhead{}| 命令可快速插入细目标注，两个相邻细目之间自然留白分隔。

\section{文字强调与列表}

\textbf{粗体}、\textit{斜体}、\texttt{等宽字体}、\underline{下划线}、{\fzkt 楷体文字}、{\fzfs 仿宋文字}、{\fzht 黑体文字}——这些都是学术写作中的常用文字效果。

\begin{itemize}
    \item 无序列表第一层：模板使用紧凑间距
    \begin{itemize}
        \item 嵌套第二层：自动缩进，行间距适中
        \item 多行长条目测试——随着人工智能技术的快速发展，越来越多的研究者开始关注模型的可解释性问题，这对于医疗、金融等高风险决策场景尤为重要。
    \end{itemize}
    \item 返回第一层
\end{itemize}

\begin{enumerate}
    \item 有序列表第一步：初始化环境
    \item 有序列表第二步：加载数据
    \begin{enumerate}
        \item 子步骤 A：数据清洗
        \item 子步骤 B：特征提取
    \end{enumerate}
    \item 有序列表第三步：模型训练
\end{enumerate}

% ============================================================
\chapter{数学公式排版}

报告模板继承 \texttt{harryopo-base.sty} 的全部数学排版能力，使用 \texttt{unicode-math} + XITS Math 引擎。

\section{行内与行间公式}

行内公式如正态分布 $X \sim \mathcal{N}(\mu, \sigma^{2})$ 与周期信号 $f(t) = A\sin(2\pi ft + \phi)$ 自然嵌入正文。

单行行间公式（equation 环境）：
\begin{equation}
\mathcal{F}(\omega) = \int_{-\infty}^{\infty} f(t)\,e^{-i\omega t}\,dt
\label{eq:fourier}
\end{equation}

\section{多行对齐}

\begin{align}
\frac{\partial \mathcal{L}}{\partial \mathbf{W}} &= \frac{1}{N} \sum_{i=1}^{N} \frac{\partial \ell(\hat{y}_i, y_i)}{\partial \mathbf{W}} \label{eq:grad_w}\\
\frac{\partial \mathcal{L}}{\partial \mathbf{b}} &= \frac{1}{N} \sum_{i=1}^{N} \frac{\partial \ell(\hat{y}_i, y_i)}{\partial \mathbf{b}} \label{eq:grad_b}
\end{align}

\section{分段函数}

\begin{equation}
\operatorname{Softmax}(\mathbf{z})_i = \frac{e^{z_i}}{\sum_{j=1}^{K} e^{z_j}}, \quad i = 1, 2, \dots, K
\label{eq:softmax}
\end{equation}

\section{定理环境}

\begin{definition}[贝叶斯公式]
对于事件 $A$ 和 $B$，若 $P(B) > 0$，则：
\begin{equation}
P(A \mid B) = \frac{P(B \mid A) \cdot P(A)}{P(B)}
\end{equation}
\end{definition}

\begin{theorem}[中心极限定理]
设 $X_1, X_2, \dots, X_n$ 为独立同分布的随机变量，均值为 $\mu$，方差为 $\sigma^{2} < \infty$，则：
\begin{equation}
\sqrt{n}\left(\frac{\bar{X}_n - \mu}{\sigma}\right) \xrightarrow{d} \mathcal{N}(0, 1)
\end{equation}
\end{theorem}

% ============================================================
\chapter{表格排版}

\section{三线表}

表\ref{tab:compare}展示了标准的三线表格式。

\begin{table}[htbp]
\centering
\begin{tabular}{lccc}
\toprule
\textbf{函数} & \textbf{公式} & \textbf{值域} & \textbf{可导性} \\
\midrule
Sigmoid & $\sigma(x) = \frac{1}{1+e^{-x}}$ & $(0, 1)$ & 全局可导 \\
Tanh    & $\tanh(x) = \frac{e^{x}-e^{-x}}{e^{x}+e^{-x}}$ & $(-1, 1)$ & 全局可导 \\
ReLU    & $\max(0, x)$ & $[0, \infty)$ & 除 $x=0$ 外可导 \\
GELU    & $x \cdot \Phi(x)$ & $(-\infty, \infty)$ & 全局可导 \\
\bottomrule
\end{tabular}
\caption{常见激活函数对比}
\label{tab:compare}
\end{table}

\section{自适应列宽表}

\begin{table}[htbp]
\centering
\begin{tabularx}{\textwidth}{lXX}
\toprule
\textbf{框架} & \textbf{优势} & \textbf{适用场景} \\
\midrule
PyTorch    & 动态计算图，调试友好，社区活跃 & 学术研究、快速原型开发 \\
TensorFlow & 生产部署成熟，Serving 生态完善 & 工业级应用、移动端部署 \\
JAX        & 函数式编程，自动向量化，高性能 & 前沿算法研究、科学计算 \\
\bottomrule
\end{tabularx}
\caption{深度学习框架选型建议}
\label{tab:framework}
\end{table}

\section{多行合并表}

\begin{table}[htbp]
\centering
\begin{tabular}{lccc}
\toprule
\multirow{2}{*}{\textbf{模型}} & \multirow{2}{*}{\textbf{参数量}} & \multicolumn{2}{c}{\textbf{准确率（\%）}} \\
\cmidrule(lr){3-4}
   &  & \textbf{验证集} & \textbf{测试集} \\
\midrule
ResNet-18  & 11.7M & 92.3 & 91.8 \\
ResNet-50  & 25.6M & 94.1 & 93.6 \\
ViT-B/16   & 86.6M & 95.8 & 95.2 \\
\bottomrule
\end{tabular}
\caption{实验配置（含多行合并）}
\label{tab:config}
\end{table}

% ============================================================
\chapter{图片插入}

\section{单图}

图\ref{fig:pipeline}使用 TikZ 内联绘制数据流水线示意。

\begin{figure}[htbp]
\centering
\begin{tikzpicture}[scale=1.1,
    box/.style={draw,thick,fill=MainColor!10,rounded corners,minimum width=2cm,minimum height=1cm,font=\small\bfseries},
    arrow/.style={->,thick,SubColor}]
    \node[box] (A) at (0,0)    {数据采集};
    \node[box] (B) at (3,0)    {数据清洗};
    \node[box] (C) at (6,0)    {特征提取};
    \node[box] (D) at (9,0)    {模型训练};
    \draw[arrow] (A) -- (B);
    \draw[arrow] (B) -- (C);
    \draw[arrow] (C) -- (D);
    \node[font=\footnotesize,below=0.6cm of B] {\fzkt 机器学习数据流水线};
\end{tikzpicture}
\caption{机器学习数据流水线（TikZ 内联绘制）}
\label{fig:pipeline}
\end{figure}

\section{子图}

\begin{figure}[htbp]
\centering
\begin{subfigure}{0.45\textwidth}
\centering
\begin{tikzpicture}[scale=0.8]
    \draw[->,blue!50,thick] (0,0) -- (2,1.5) node[right,font=\footnotesize] {$y=x$};
    \fill[blue!30] (1,1) circle (0.05);
    \fill[blue!30] (0.8,1.2) circle (0.05);
    \fill[blue!30] (1.2,0.9) circle (0.05);
    \node[font=\footnotesize] at (1,-0.5) {(a) 线性关系};
\end{tikzpicture}
\caption{线性关系}
\end{subfigure}
\hfill
\begin{subfigure}{0.45\textwidth}
\centering
\begin{tikzpicture}[scale=0.8]
    \draw[->,red!50,thick] (0,0) .. controls (0.8,1.2) and (1.2,0.3) .. (2,1.5);
    \fill[red!30] (0.5,0.8) circle (0.05);
    \fill[red!30] (1,0.6) circle (0.05);
    \fill[red!30] (1.5,1.1) circle (0.05);
    \node[font=\footnotesize] at (1,-0.5) {(b) 非线性关系};
\end{tikzpicture}
\caption{非线性关系}
\end{subfigure}
\caption{数据分布的子图对比}
\label{fig:subfig}
\end{figure}

% ============================================================
\chapter{代码高亮}

\section{Python 代码}

模板预置 Python 代码高亮风格（\texttt{pystyle}）。

\begin{lstlisting}[style=pystyle, caption={数据预处理示例}, label=lst:preprocess]
import numpy as np
from sklearn.preprocessing import StandardScaler

def preprocess(X_train, X_test):
    """标准化预处理"""
    scaler = StandardScaler()
    X_train_scaled = scaler.fit_transform(X_train)
    X_test_scaled  = scaler.transform(X_test)
    return X_train_scaled, X_test_scaled

# 加载数据
X, y = load_dataset("cifar10")
X_norm = (X - np.mean(X)) / np.std(X)
\end{lstlisting}

\section{Shell 命令}

\begin{lstlisting}[style=bashstyle, caption={环境配置}, label=lst:env]
# 创建虚拟环境
python -m venv venv && source venv/bin/activate

# 安装依赖
pip install torch==2.1.0 transformers datasets tiktoken

# 运行训练
python train.py --model llama2 --epochs 3 --batch_size 32 \
    --learning_rate 2e-5 --warmup_steps 500
\end{lstlisting}

% ============================================================
\chapter{文献引用与交叉引用}

\section{常规引用}

学术写作离不开规范的文献引用。以下为常用引用的展示：

\begin{itemize}
    \item 深度学习经典综述：\cite{lecun2015deep}
    \item ResNet 残差网络：\cite{he2016deep}
    \item Transformer 原始论文：\cite{vaswani2017attention}
    \item 神经网络正则化：\cite{srivastava2014dropout}
    \item BN 加速训练：\cite{ioffe2015batch}
    \item 中文文献范例：\cite{zhou2016machine}
\end{itemize}

\section{交叉引用}

交叉引用命令 \verb|\ref{}| 和 \verb|\eqref{}| 可自动编号：文字排版在第一章；公式\eqref{eq:fourier}为傅里叶变换；表\ref{tab:compare}为激活函数对比；图\ref{fig:pipeline}为数据流水线；各节代码如\ref{lst:preprocess}和\ref{lst:env}。

% ============================================================
\chapter{模板选项速查}

\section{文档类参数}

\begin{itemize}
    \item \verb|green| — 切换绿色主题
    \item \verb|dark| — 切换深紫色主题
\end{itemize}

\section{核心命令一览}

\begin{table}[htbp]
\centering
\begin{tabular}{lp{7cm}}
\toprule
\textbf{命令} & \textbf{功能} \\
\midrule
\verb|\maketitle| & 生成封面页 \\
\verb|\subtitle{...}| & 副标题 \\
\verb|\institute{...}| & 机构名称 \\
\verb|\introduction| & 生成引言章节 \\
\verb|\subhead{标题}| & 括号细目标题 \\
\verb|\printfunding| & 打印基金信息 \\
\verb|\fzkt| / \verb|\fzfs| / \verb|\fzht| & 楷体 / 仿宋 / 黑体 \\
\bottomrule
\end{tabular}
\caption{报告模板核心命令速查}
\end{table}

% --- 参考文献 ---
\begin{thebibliography}{9}

\bibitem{lecun2015deep} LeCun Y, Bengio Y, Hinton G. Deep learning[J]. \emph{Nature}, 2015, 521(7553): 436--444.

\bibitem{he2016deep} He K, Zhang X, Ren S, et al. Deep residual learning for image recognition[C]. \emph{CVPR}, 2016: 770--778.

\bibitem{vaswani2017attention} Vaswani A, Shazeer N, Parmar N, et al. Attention is all you need[C]. \emph{NeurIPS}, 2017: 5998--6008.

\bibitem{srivastava2014dropout} Srivastava N, Hinton G, Krizhevsky A, et al. Dropout: A simple way to prevent neural networks from overfitting[J]. \emph{JMLR}, 2014, 15(1): 1929--1958.

\bibitem{ioffe2015batch} Ioffe S, Szegedy C. Batch normalization: Accelerating deep network training by reducing internal covariate shift[C]. \emph{ICML}, 2015: 448--456.

\bibitem{deng2009imagenet} Deng J, Dong W, Socher R, et al. ImageNet: A large-scale hierarchical image database[C]. \emph{CVPR}, 2009: 248--255.

\bibitem{zhou2016machine} 周志华. 机器学习[M]. 北京: 清华大学出版社, 2016.

\end{thebibliography}

\end{document}
